\documentclass[10pt,a4paper]{amsart}
\usepackage[margin=19mm]{geometry}
\usepackage{amsmath,amssymb,mathtools}
\usepackage[scr=rsfs,cal=euler]{mathalfa}
\usepackage{xfrac}
\usepackage[unicode,hidelinks,bookmarks=false]{hyperref}
\newcommand{\ie}{i.e.\ }
\newcommand{\cf}{cf.\ }
\newcommand{\resp}{respectively\ }
\newcommand{\Subsection}{\subsection}
\providecommand{\nobreakdash}{\nobreak\hbox{-}}
\setlength{\emergencystretch}{2em}
\multlinegap0pt
\usepackage{bm}
\usepackage{bbm}
\usepackage{dsfont}
\usepackage{xspace}
\usepackage{cancel}
\usepackage{mathtools}
\usepackage{amsthm}
\makeatletter
\@ifundefined{theorem}{\newtheorem{theorem}{Theorem}}{}
\@ifundefined{prop}{\newtheorem{prop}[theorem]{Proposition}}{}
\makeatother
\newcommand{\eight}{{\mathbf 8}}
\newcommand{\lamm}{\Lambda^-}
\newcommand{\qpx}{\Omega_{240}}
\title[A partial order on the 240 packings of PG(3,2)]{A partial order on the 240 packings of PG(3,2)}
\author{R.M. Green}
\date{Source: arXiv:2511.03040v1 (CC BY 4.0)}
\begin{document}
\maketitle

\noindent\emph{Source and licence.} A partial order on the 240 packings of PG(3,2). Version arXiv:2511.03040v1 (CC BY 4.0), distributed under the Creative Commons Attribution 4.0 International licence, as recorded in the arXiv licence metadata for this exact version and shown on the abstract page https://arxiv.org/abs/2511.03040v1. Full rights notice: licenses/arxiv-2511.03040-CC-BY-4.0.txt.\par\smallskip

\noindent\textit{Editorial scope.} Counted unit: the Prop whose source label is \texttt{prop:signact} in the article (source0.tex lines 1072--1085, source label \texttt{prop:signact}), delivered together with the article's own complete original proof of it (source0.tex lines 1087--1150, one \texttt{proof} environment of 64 lines). No mathematics is written, repaired or re-derived: statement and proof are the article's own text line for line. Required context retained uncounted: none. Every definition, display and label of those lines is retained unchanged. Omitted from this package and not redistributed: 1--1071, 1086--1086, 1151--2070. Nothing inside a retained excerpt is omitted or altered: every retained body is delivered exactly as the article's lines are written, so no omission receipt of any kind accompanies this package. Citations: the retained excerpts carry no citation command at all, so no bibliography entry of the article is transcribed and no reference is redistributed. Reference adaptation: none was needed and none was made; every reference inside the delivered unit resolves inside it, so no unresolved reference or citation is produced. Printed slips kept literal: no printed word, symbol, hypothesis or number of the counted statement or of its proof was altered. Layout and class: the article's own class is \texttt{amsart} with packages including babel, amsmath, graphicx, hyperref, amsthm, amssymb, amsfonts, amscd, fullpage, subfig, tikz, xspace, adjustbox; the wrapper is the lane's pinned \texttt{amsart} editorial wrapper at 10pt on A4, which restates the theorem-like environments the retained text needs and adds the article's own macro definitions that the retained text needs (\texttt{\textbackslash eight}, \texttt{\textbackslash lamm}, \texttt{\textbackslash qpx}), with extra wrapper packages (bm, bbm, dsfont, xspace, cancel, mathtools). The delivered numbering is the unit's own. Assets: no retained span contains a figure environment, a graphics-inclusion command, a PDF-page inclusion, a table or a TikZ picture; no image file is copied into this package, the package declares no asset dependency and no image file is redistributed with it. Licence: version arXiv:2511.03040v1 of this article is distributed under the Creative Commons Attribution 4.0 International licence, as recorded in the arXiv licence metadata for this exact version and shown on the abstract page https://arxiv.org/abs/2511.03040v1. A full rights notice is delivered at licenses/arxiv-2511.03040-CC-BY-4.0.txt. Characters: every retained excerpt is pure ASCII, so the delivered engine, XeLaTeX with its default Latin Modern text font, reports no missing character.
\begin{prop}\label{prop:signact}
Let $x \in \qpx$ be a packing of $PG(3, 2)$ identified with an $\eight$-labelled Fano plane. Let $l \in \eight$ 
be the basepoint of $x$, let $e$ be one of the labels in $x$, and denote the three triples containing
$e$ by $\{e, a_1, b_1\}$, $\{e, a_2, b_2\}$, and $\{e, a_3, b_3\}$. 
\begin{itemize}
\item[{\rm (i)}]{We have $\overline{(l, e)}x = (a_1, b_1)x$ and $\overline{(a_1, b_1)}x = (l, e)x$.}
\item[{\rm (ii)}]{We have $(a_1, b_1)x = (a_2, b_2)x = (a_3, b_3)x$ and
$\overline{(a_1, b_1)}x = \overline{(a_2, b_2)}x = \overline{(a_3, b_3)}x$.}
\item[{\rm (iii)}]{For any $1 \leq i < j \leq 8$, there exist $1 \leq c < d \leq 8$ such that
$\overline{(i, j)}x = (c, d)x$.}
\item[{\rm (iv)}]{For any $1 \leq i < j \leq 8$, there exist $1 \leq c < d \leq 5$ such that either
$(i, j)x = (c, d)x$ or $(i, j)x = \overline{(c, d)}x$.}
\end{itemize}
\end{prop}

\begin{proof}
By symmetry, it suffices to verify (i) and (ii) for a specific choice of $x$, $l$, and $e$. We will consider 
the case $l = 8$, $x = x_0$, and $e = 1$.

We have reduced (i) to the case $\overline{(1, 8)}x_0 = (2, 7)x_0$ and $\overline{(2, 7)}x_0 = (1, 8)x_0$.
The action of $\overline{(1, 8)}$ on $\lamm$ sends $$
\{ v_{127},
v_{136},
v_{145},
v_{235},
v_{246},
v_{347},
v_{567},
v_{8}\} \quad \text{to} \quad 
\{ v_{127},
v_{136},
v_{145},
v_{467},
v_{357},
v_{256},
v_{234},
v_{8}\}
,$$ where we have listed the eight components in the corresponding order. By inspection, the action of $(2, 7)$ on
$x_0$ has the same effect. Similarly, the action of $\overline{(2, 7)}$ on $\lamm$ sends $$
\{ v_{127},
v_{136},
v_{145},
v_{235},
v_{246},
v_{347},
v_{567},
v_{8}\} \quad \text{to} \quad
\{ v_{1},
v_{458},
v_{368},
v_{235},
v_{246},
v_{347},
v_{567},
v_{278}\}
,$$ which by inspection agrees with the action of $(1, 8)$ on $x_0$. This completes the proof of (i).

The reason that $(a_1, b_1)$, $(a_2, b_2)$, and $(a_3, b_3)$ have the same effect on $x$ in (ii) is that they
all have the same effect as the action of $\overline{(l, e)}$ on $x$, by (i). Similarly,
$\overline{(a_1, b_1)}$, $\overline{(a_2, b_2)}$, and $\overline{(a_3, b_3)}$ all have the same effect on $x$
as $(l, e)$, which proves (ii).

If we have $l \in \{i, j\}$ then (iii) follows from the first assertion of (i); if not, then (iii) follows from
the second assertion of (i).

To prove (iv), let us first suppose that $l \in \{i, j\}$. Define $e$ so that $\{l, e\} = \{i, j\}$.
Part (iii) provides three signed transpositions, $\overline{(a_1, b_1)}$, $\overline{(a_2, b_2)}$, and 
$\overline{(a_3, b_3)}$, that have the same effect on $x$ as $(l, e)$. The set $\{l, e, a_1, b_1, a_2, b_2, a_3, b_3\}$
has size $8$, so at least one of the transpositions $(l, e)$, $\overline{(a_1, b_1)}$, $\overline{(a_2, b_2)}$, and 
$\overline{(a_3, b_3)}$ is of the form $(c, d)$ or $\overline{(c, d)}$ such that $\{c, d\} \cap \{6, 7, 8\} = \emptyset$, 
as required.

It remains to deal with the case where $l \not\in \{i, j\}$. Let $e$ be the label for which $\{e, i, j\}$
is a triple of $x$, and let $\{e, a, b\}$ and $\{e, a', b'\}$ be the other triples that contain $e$.
If at least one of the pairs $\{i, j\}$, $\{a, b\}$ and $\{a', b'\}$ is disjoint from $\{6, 7, 8\}$, we define
$\{c, d\}$ to be this pair, and it follows from (iii) that $(c, d)$ and $(i, j)$ have the same effect on $x$,
proving (iv). Otherwise, the pair $\{l, e\}$ is disjoint from $\{6, 7, 8\}$, and we define $\{c, d\} = \{l, e\}$.
It follows from (i) that $\overline{(c, d)}$ and $(i, j)$ have the same effect on $x$, and this completes the proof.
\end{proof}

\end{document}
